% \section{Layer Summary}

% Building on the prior-work positioning above,
% this dissertation organizes its contributions by where 
% a defense intervenes in an emerging software system.
% Some defenses constrain behavior at the moment of high-stakes execution.
% Others reveal or reduce harmful behavior induced earlier by 
% opaque models or weakly trusted information sources.
% Under this view, the four works form the following 
% connected defense layers.

% \noindent \textbf{Layer I: Demystifying invariant effectiveness for securing smart contracts.}
% it proposes runtime invariant generation  and enforcement techniques
% for smart contract security, where invariants are rules that 
% should continue to hold during normal execution.

% \noindent \textbf{Layer II: Enforcing control flow integrity on DeFi smart contracts.}
% it proposes algorithms for control flow integrity 
% enforcement across contract interactions, 
% with minimal runtime overhead; informally, 
% this layer ensures that contract functions are called 
% in legitimate patterns.

% \noindent \textbf{Layer III: Auditing malicious scam endpoints in production LLMs.}
% it audits LLMs when generating code under realistic developer-style prompts, 
% and finds that LLMs can generate malicious scam URLs with high frequency. 
% This layer reveals a new attack vector in LLMs, and proposes 
% a practical detection system to audit malicious endpoints in the pre-training 
% data of LLMs.

% \noindent \textbf{Layer IV: Defending Search- and Documentation-Level Poisoning of Coding Agents.}
% it studies how poisoned external content retrieved by coding agents 
% can propagate into harmful coding outcomes, and how to defend against this pipeline.

